%%%%%%%%%%  scan idx 310  |  folio 297  %%%%%%%%%%
\SGAPrintPage{297}%
% Exposé V: back to the volume's parf
\renewcommand{\parf}{\operatorname{parf}}
\sgaStableTarget{sga6.e05.expose}\section*{\underline{Exposé V}\\[3pt] \underline{\underline{GENERALITIES ON \(\lambda\)-RINGS}}}
\begin{center}
\underline{by P. BERTHELOT}
\end{center}

The theory of \(\lambda\)-rings is useful whenever one uses rings such as the ring of classes of vector bundles on a ringed site, the representation ring of a group, etc..., for which the exterior-power operations define maps \(\lambda^i\). We shall use it extensively in the remainder of the Seminar, so it is necessary first to develop its general formalism. This exposé is thus entirely independent of the rest of the seminar and can serve users of \(\lambda\)-rings outside algebraic geometry.

The first section essentially gives Hirzebruch's theory of ``multiplicative sequences of polynomials'' (cf. [2]); the functorial viewpoint adopted here has the advantage, besides its generality, of considerably reducing the calculations. The following sections will show how, to every augmented \(\lambda\)-ring, one can associate a ``Chern ring'' that permits the development of a theory of Chern classes. This is the theory that will enable us to state the Riemann--Roch theorem in the general form of \hyperlink{sga6.e08.expose}{Exp. VIII}.

\sgaStableTarget{sga6.e05.sec.1}\section*{1. Universal polynomials}

\sgaStableTarget{sga6.e05.para.1-1}1.1. Let \(K\) be a commutative ring with unit, fixed throughout this section. We denote by \(\mathcal C\) the category of filtered commutative \(K\)-algebras with unit, whose filtration is decreasing and such that the group of elements of filtration zero is the whole algebra. We define a functor on \(\mathcal C\), with values in the category of groups
%%%%%%%%%%  scan idx 311  |  folio 298  %%%%%%%%%%
\SGAPrintPage{298}%
that are abelian, as follows: for \(A\in\operatorname{Ob}(\mathcal C)\), we associate to \(A\) the set of formal series \(\psi\in A[[t]]\) satisfying the following conditions: if \(\psi=a_0+a_1t+\ldots+a_nt^n+\ldots\), then
\sgaStableTarget{sga6.e05.eq.1-1-1}\begin{equation}\tag{1.1.1}
  a_0=1;\qquad \forall i\geq 1,\quad a_i\in\Fil^i A .
\end{equation}
This set will be denoted by \(\Ghat_K(A)\), or by \(\Ghat(A)\) when no confusion is possible; it is clear that it is an abelian group under multiplication of formal series, and that this does define a functor
\[
  \Ghat_K:\underline{K\text{-Alg.fil.}}\longrightarrow \underline{\mathrm{Ab}}.
\]
For \(A\) endowed with the coarse filtration, \(\Ghat(A)\) is the group of formal series of augmentation \(1\), which will also be denoted by \(1+A[[t]]^+\).

\sgaStableTarget{sga6.e05.lem.1-2}\textbf{Lemma 1.2.} \uline{The functor} \(\Ghat_K\) \uline{is representable}.

Let \((X_i)_{i\in N,\ i\geq 1}\) be a family of indeterminates; we assign the weight \(i\) to \(X_i\), and filter the ring \(K[X_i]_{i\geq 1}\) by weight with respect to the \(X_i\). We then define a bijection:
\[
  \Ghat_K(A)\xrightarrow{\sim}
  \Hom_{K\text{-Alg.fil.}}\bigl(K[X_i]_{i\geq 1},A\bigr)
\]
by associating to every series \(\psi=1+a_1t+\ldots+a_nt^n+\ldots\in\Ghat_K(A)\) the homomorphism \(f_\psi\) determined by \(f_\psi(X_i)=a_i\). This homomorphism is compatible with the filtrations, and we thus obtain an isomorphism of functors.

Put \(\phi=1+X_1t+\ldots+X_nt^n+\ldots\); for every filtered \(K\)-algebra \(A\) and every series \(\psi\in\Ghat(A)\), there exists exactly one homomorphism
%%%%%%%%%%  scan idx 312  |  folio 299  %%%%%%%%%%
\SGAPrintPage{299}%
\(f_\psi\) from \(K[X_i]_{i\geq 1}\) to \(A\) such that \(\Ghat(f_\psi)(\phi)=\psi\).

\sgaStableTarget{sga6.e05.para.1-3}1.3. For every \(A\in\operatorname{Ob}(\mathcal C)\) we define a filtration on \(\Ghat(A)\) as follows: a series \(\psi=1+a_1t+\ldots\) will be said to be of filtration \(\geq k\) if \(a_i=0\) for \(1\leq i<k\). We denote the subgroup so obtained by \(F^k\Ghat(A)\); this is a functor on \(\mathcal C\), and from now on we shall regard \(\Ghat\) as a functor with values in the category \(\underline{\mathrm{Ab}}.\underline{\mathrm{fil}}.\) of filtered abelian groups by means of this filtration (with \(F^0\Ghat=\Ghat\)).

Let \(H\) be a functor on \(\mathcal C\) with values in \(\underline{\mathrm{Ab}}\,\underline{\mathrm{fil}}\); we shall assume that \(H\) satisfies the following two conditions:
\[
\begin{array}{ll}
\text{i)} & \forall A\in\operatorname{Ob}(\mathcal C),\ H(A)\text{ is separated and complete}.\\[0.3em]
\text{ii)} & \text{If }A\in\operatorname{Ob}(\mathcal C)\text{ is a graded }K\text{-algebra, endowed with the filtration}\\
& \text{associated with its grading, and if }G\text{ is a finite group acting on }A\\
& \text{by graded morphisms, then }H(A^G)\xrightarrow{\sim} H(A)^G\\
& \text{(the exponent }G\text{ denoting invariants under }G\text{)}
\end{array}
\]

In fact, we shall need these conditions only when \(A\) is a polynomial ring and \(G\) is a finite product of symmetric groups acting by permutation of the indeterminates.

\sgaStableTarget{sga6.e05.ex.unnamed}\textbf{Examples}: 1)~\(H(A)=\Ghat(A)\). \sgaStableTarget{sga6.e05.ex.1-3-2}2)~\(H(A)=1+A[[t]]^+\); this functor will be denoted by \(\Gg\) (\(g\), the initial of ``grossière''). 3)~\(H(A)=\Ga(A)\), where \(\Ga(A)\) denotes the group, under addition, of formal series such that \(a_0=0\), \(a_i\in\Fil^i A\) for \(i\geq 1\), endowed with the filtration analogous to that of \(\Ghat\).

\sgaStableTarget{sga6.e05.para.1-4}1.4. For every \(n\), we shall denote by \(f_n\) the endomorphism of \(K[X_1,\ldots,X_n]\)
%%%%%%%%%%  scan idx 313  |  folio 300  %%%%%%%%%%
\SGAPrintPage{300}%
defined by \(f_n(X_i)=X_i\) for \(i<n\), \(f_n(X_n)=0\). It factors through the inclusion of \(K[X_1,\ldots,X_{n-1}]\) in \(K[X_1,\ldots,X_n]\), which shows that the homomorphism from \(H(K[X_1,\ldots,X_{n-1}])\) to \(H(K[X_1,\ldots,X_n])\) is injective, and will allow us to identify the former with a subgroup of the latter; similarly, we shall regard \(H(K[X_1,\ldots,X_n])\) as contained in \(H(K[X_i]_{i\geq 1})\).

Notice that \(f_1\) is precisely the augmentation of \(K[X]\).

\sgaStableTarget{sga6.e05.para.1-5}1.5. We embed \(K[X_1,\ldots,X_n]\) in \(K[T_1,\ldots,T_n]\) by sending \(X_i\) to the \(i\)-th symmetric function in the \(T_j\); letting the symmetric group \(\mathfrak S_n\) act by permutation of the \(T_j\), and using the theorem on symmetric polynomials and (1.3.ii)), we see that the resulting homomorphism from \(H(K[X_1,\ldots,X_n])\) to \(H(K[T_1,\ldots,T_n])\) is injective, which will again allow us to identify the former with a subgroup of the latter:
\[
  H(K[X_1,\ldots,X_n])\simeq H(K[T_1,\ldots,T_n]^{\mathfrak S_n})\xrightarrow{\sim} H(K[T_1,\ldots,T_n])^{\mathfrak S_n}\subset H(K[T_1,\ldots,T_n])
\]
We denote by
\[
  u_i^n:K[X]\longrightarrow K[T_1,\ldots,T_n]
\]
the homomorphism defined by \(u_i^n(X)=T_i\), and put
\[
  u_H^n=\sum_{i=1}^n H(u_i^n):H(K[X])\longrightarrow H(K[T_1,\ldots,T_n]).
\]
The symmetric group permutes the morphisms \(H(u_i^n)\), and therefore leaves the morphism \(u_H^n\) invariant, which consequently takes its values in
%%%%%%%%%%  scan idx 314  |  folio 301  %%%%%%%%%%
\SGAPrintPage{301}%
\(H(K[X_1,\ldots,X_n])\) by (1.3.ii)).

Finally, observe that for \(H=\Ghat\), \(u_H^n(1+Xt)=\prod\limits_{i=1}^n(1+T_it)=1+X_1t+\ldots+X_nt^n\).

\sgaStableTarget{sga6.e05.thm.1-6}\textbf{Theorem 1.6.} \uline{The map from} \(\Hom_{\underline{\mathrm{Ab}}\,\underline{\mathrm{fil}}}(\Ghat_K,H)\) \uline{to} \(H(K[X])\) \uline{which}, \uline{to a morphism of functors}, \uline{assigns its value on the element} \(1+Xt\) \uline{of} \(\Ghat_K(K[X])\), \uline{is injective}; \uline{its image is the subgroup of} \(H(K[X])\) \uline{formed by the elements} \(\theta\) \uline{such that} (\uline{with the notation of} 1.4. \uline{for the} \(f_n\)):
\[
\begin{array}{ll}
\text{i)}& H(f_1)(\theta)=0.\\[0.3em]
\text{ii)}& \forall n\geq1,\quad
H(f_n)\bigl(u_H^n(\theta)\bigr)\equiv u_H^n(\theta)
\ \mathrm{mod.}\ F^nH(K[X_1,\ldots,X_n]).
\end{array}
\]

Let us first prove injectivity. By Lemma 1.2, it is enough to see that a morphism of functors \(\Phi\) vanishing on \(1+Xt\) vanishes on \(\varphi=1+X_1t+\ldots+X_nt^n+\ldots\). The commutative diagram:
\sgaStableTarget{sga6.e05.eq.1-6-1}\[
\begin{tikzcd}
\Ghat(K[X]) \arrow[r,"u^n_{\Ghat}"] \arrow[d,"\Phi"'] &
\Ghat(K[X_1,\ldots,X_n]) \arrow[d,"\Phi"]\\
H(K[X]) \arrow[r,"u_H^n"] &
H(K[X_1,\ldots,X_n])
\end{tikzcd}
\tag{1.6.1}
\]
shows that for every \(n\), \(\Phi(1+X_1t+\ldots+X_nt^n)=0\).

Now \(\varphi\equiv 1+X_1t+\ldots+X_nt^n\ \mathrm{mod.}\ F^{n+1}\Ghat\); hence \(\Phi(\varphi)\equiv \Phi(1+\ldots+X_nt^n)\ \mathrm{mod.}\ F^{n+1}H\), and consequently \(\Phi(\varphi)\equiv0\ \mathrm{mod.}\ F^{n+1}H\), for every \(n\); since \(H(K[X_i]_{i\geq 1})\) is separated and complete, \(\Phi(\varphi)=0\).

Now let \(\Phi\) be any element of \(\Hom(\Ghat,H)\), and put
